\section{The Physics Validation: The \texorpdfstring{$E_8$}{E8} Substrate and the Geometric Cascade}
\label{sec:physics_validation}

As derived in \textit{Persistence Mechanics: A Universal Architecture of Persistence} \cite{meyer_informational_2026}, the vacuum is the unique resource-constrained computational substrate satisfying the six architectural pillars of persistence: an $E_8$ lattice projected to $D=4$ spacetime, with each localized node possessing strictly finite channel capacity ($\nu = 16$). This substrate contains four key invariants: the temporal resonance $\Delta=43$, the chiral capacity $\nu=16$, the interaction rank $\sigma=5$, and the topological boundary $\chi=2$, all derived from the $E_8 \rightarrow D_4 \oplus D_4$ projection.

The companion work \textit{Persistence Mechanics: Entropic Action} \cite{meyer_ieaction_2026} proposes the continuum functional operator forms; this paper examines numerical allocation laws. The general results of \cref{sec:nested_persistence} constrain admissible resource use, reversible encoding, and interface transfer. Applying them to the vacuum additionally requires its workload, operational encoding, response law, and selection objective. The formulas below are proposed substrate-specific allocation laws whose consequences can be checked independently. Global solvency alone does not prove that these allocations are unique; proving that the substrate selects their additional premises remains a separate obligation.

If the spatial geometry, temporal updates, and internal force generators all attempt to saturate the node simultaneously, local complexity exceeds the 16 available channels and the causal structure becomes unresolvable. Because the vacuum persists, its structure must satisfy the three rules of Nested Persistence (\cref{sec:nested_persistence}). The $\nu = 16$ budget is therefore partitioned across orthogonal interaction sectors:

\textbf{Rule 1 (Capacity Partitioning)} allocates the $\nu = 16$ chiral capacity across orthogonal geometric sectors. This yields \textbf{System 4: The Gauge Partition}: the saturation ratio (Strong Coupling $\alpha_s$), the spacetime split (Weak Mixing Angle $\sin^2\theta_W$), and the irreducible geometric frustration of projecting 5-fold internal symmetry onto 4-fold spacetime (Jarlskog Invariant $J$). Because the temporal and color sectors are geometrically isolated by this partition, the Strong CP problem is structurally prohibited without requiring an axion.

\textbf{Rule 3 (Impedance Matching)} stabilizes the boundary between the high-frequency vacuum resonance and local topological defects. This yields \textbf{System 5: The Resonant Interface}, generating the electroweak mass scale: the Higgs VEV ($v$), Fermi constant ($G_F$), and the minimum resolvable mass (Electron Yukawa $y_e$), the floor below which vacuum fluctuations erase distinguishable states.

\textbf{Rules 1 \& 2 (Capacity Partitioning and Reversible Encoding)} govern the residual bandwidth, now attenuated across the geometric bulk. This yields \textbf{System 6: The Geometric Depth}, identified as Gravity. The attenuation depth derives the Planck scale ($M_P$); the irreducible bulk resolution floor derives the Vacuum Energy Density ($\rho_{vac}$), resolving both the Gravitational Hierarchy and the Vacuum Catastrophe.

\subsection{Note on Absolute Scale}

A discrete lattice inherently produces pure dimensionless ratios such as $v/m_e \approx 4.8 \times 10^5$ and $M_P/m_e \approx 2.4 \times 10^{22}$; the electron Yukawa then follows as $y_e=\sqrt{2}\,m_e/v$. These ratios are the framework's proposed predictions, conditional on the operational premises stated with each one. To express them as absolute values in laboratory units (MeV, GeV), one scale must be fixed --- exactly as expressing the speed of light in m/s requires the definition of a meter. This anchoring is intrinsic to any scale-invariant framework and introduces no free parameter.

We adopt $m_e \approx 0.511\,\mathrm{MeV}$ as the \textbf{unit anchor}. The proposed interpretation is that the electron is the lightest stable charged defect, the Margin floor of System 5. That identification is a spectral hypothesis; the encoding theorems bound metadata but do not by themselves produce a mass spectrum (\cref{sec:Synthesis}).

The proposed higher scales ($v$, $M_H$, $M_P$) are geometric multiples of this anchor; the derivations supply dimensionless ratios and the anchor converts them to laboratory units. No coefficient is fitted to data, but each ratio depends on stated workload, response and matching premises, and the absolute W and Z masses additionally require the interface response discussed in \cref{subsec:weak_mixing}. The span runs from $\alpha_G\sim10^{-45}$ to $M_P\sim10^{19}$~GeV.